\documentclass[10pt,a4paper]{amsart}
\usepackage[margin=19mm]{geometry}
\usepackage{amsmath,amssymb,mathtools}
\usepackage[scr=rsfs,cal=euler]{mathalfa}
\usepackage{xfrac}
\usepackage[unicode,hidelinks,bookmarks=false]{hyperref}
\newcommand{\ie}{i.e.\ }
\newcommand{\cf}{cf.\ }
\newcommand{\resp}{respectively\ }
\newcommand{\Subsection}{\subsection}
\providecommand{\nobreakdash}{\nobreak\hbox{-}}
\setlength{\emergencystretch}{2em}
\multlinegap0pt
\usepackage{bm}
\usepackage{bbm}
\usepackage{dsfont}
\usepackage{xspace}
\usepackage{cancel}
\usepackage{mathtools}
\usepackage{cleveref}
\usepackage[shortlabels]{enumitem}
\usepackage{amsthm}
\makeatletter
\@ifundefined{lemma}{\newtheorem{lemma}{Lemma}}{}
\makeatother
\newcommand{\RR}{\mathbb{R}}
\newcommand{\dd}{\mathrm{d}}
\newcommand{\del}{\partial}
\title[Analytic Fourier Integral Operators and a Problem from Seism]{Analytic Fourier Integral Operators and a Problem from Seismic Inversion (Lemma 2)}
\author{Leonard Busch}
\date{Source: arXiv:2505.08472v2 (CC BY 4.0)}
\begin{document}
\maketitle

\noindent\emph{Source and licence.} Analytic Fourier Integral Operators and a Problem from Seismic Inversion. Version arXiv:2505.08472v2 (CC BY 4.0), distributed under the Creative Commons Attribution 4.0 International licence, as recorded in the arXiv licence metadata for this exact version and shown on the abstract page https://arxiv.org/abs/2505.08472v2. Full rights notice: licenses/arxiv-2505.08472-CC-BY-4.0.txt.\par\smallskip

\noindent\textit{Editorial scope.} Counted unit: the Lemma printed as Lemma 2 of the article (source0.tex lines 857--879, source label \texttt{eq:injmatrix}), delivered together with the article's own complete original proof of it (source0.tex lines 880--950, one \texttt{proof} environment of 71 lines). No mathematics is written, repaired or re-derived: statement and proof are the article's own text line for line. Required context retained uncounted: as:n1 (lines 808--819); as:n2 (lines 808--819). Every definition, display and label of those lines is retained unchanged. Omitted from this package and not redistributed: 1--807, 820--856, 951--11890. Nothing inside a retained excerpt is omitted or altered: every retained body is delivered exactly as the article's lines are written, so no omission receipt of any kind accompanies this package. Citations: the retained excerpts carry no citation command at all, so no bibliography entry of the article is transcribed and no reference is redistributed. Reference adaptation: none was needed and none was made; every reference inside the delivered unit resolves inside it, so no unresolved reference or citation is produced. Printed slips kept literal: no printed word, symbol, hypothesis or number of the counted statement or of its proof was altered. Layout and class: the article's own class is \texttt{scrartcl} with packages including templ, bm, mathrsfs, tikz, pgf, xparse; the wrapper is the lane's pinned \texttt{amsart} editorial wrapper at 10pt on A4, which restates the theorem-like environments the retained text needs and adds the article's own macro definitions that the retained text needs (\texttt{\textbackslash RR}, \texttt{\textbackslash dd}, \texttt{\textbackslash del}), with extra wrapper packages (bm, bbm, dsfont, xspace, cancel, mathtools, cleveref, enumitem). The delivered numbering is the unit's own. Assets: no retained span contains a figure environment, a graphics-inclusion command, a PDF-page inclusion, a table or a TikZ picture; no image file is copied into this package, the package declares no asset dependency and no image file is redistributed with it. Licence: version arXiv:2505.08472v2 of this article is distributed under the Creative Commons Attribution 4.0 International licence, as recorded in the arXiv licence metadata for this exact version and shown on the abstract page https://arxiv.org/abs/2505.08472v2. A full rights notice is delivered at licenses/arxiv-2505.08472-CC-BY-4.0.txt. Characters: every retained excerpt is pure ASCII, so the delivered engine, XeLaTeX with its default Latin Modern text font, reports no missing character.
	\begin{enumerate}[label=$\mathrm{ph}$\arabic*.]%\setcounter{enumii}{-2}
		\item\label{as:n1} $\Lambda$ satisfies the Bolker condition at $(\hat z, \hat\zeta,\hat x, \hat \xi)$.
		\item\label{as:n2} In an open neighborhood of $(\hat z,\hat \zeta, \hat x, \hat \xi)$, the map 
				\[
					\Lambda \ni (z,\zeta; x,\xi)\mapsto (z,x,\xi)
				\]
				is an immersion.
		\item\label{as:n3} In an open neighborhood of $(\hat z,\hat \zeta, \hat x, \hat \xi)$, we have
			\[
				\Lambda \ni (z,\zeta ;x,\xi) \implies \zeta \neq 0\,.
			\]
	\end{enumerate}

\begin{lemma}
	Let $\Lambda$ be a canonical relation associated to some non-degenerate analytic phase $\varphi$ near some $(\hat z, \hat \zeta, \hat x, \hat \xi) \in \Lambda$, and let $(\hat \zeta,\hat \xi) = (\varphi_z,-\varphi_x)(\hat z, \hat x, \hat \eta)$ for some $(\hat z, \hat x, \hat \eta) \in \Sigma_\varphi$.

	Let the Bolker condition~\ref{as:n1} be satisfied at $(\hat z,\hat \zeta;\hat x,\hat\xi)$. We have $\dd\pi_L\vert_{(\hat z, \hat \zeta, \hat x, \hat \xi)}$ is injective if and only if for all curves $(z,x,\eta)(t)$ on $\Sigma_\varphi$ with $(z,x,\eta)(0) = (\hat z, \hat x, \hat \eta)$ we have
\begin{equation}\label{eq:injmatrix}
(\dot z(0), \del_x\varphi_{z}(\hat z, \hat x, \hat \eta)\dot x(0) +  \del_\eta\varphi_{z}(\hat z, \hat x, \hat \eta)\dot\eta(0)) = 0 \implies (\dot x(0), \del_\eta\varphi_{x}(\hat z, \hat x, \hat \eta)\dot\eta(0)) = 0\,.
%\varphi_\eta(\hat z, \hat x, \hat \eta) = 0 \text{ implies that the } (N\times (n+N'))\text{-matrix}\quad \left((\del_\eta\varphi_{z})(\hat z,\hat x,\hat \eta), (\del_x\varphi_{z})(\hat z,\hat x,\hat \eta)\right) \quad\text{is injective}\,.
\end{equation}
Furthermore, if $\Lambda$ satisfies assumption~\ref{as:n2}, 
%%%then there exists a matrix $B$ so that for all curves $(z,x,\eta)(t)$ on $\Sigma_\varphi$ with $(z,x,\eta)(0) = (\hat z, \hat x, \hat \eta)$, and $\dot z(0) = 0$, for the $N\times (N'+n)$-matrix $(\del_x\varphi_z, \del_\eta\varphi_z)(\hat z,\hat x,\hat \eta)$ we have % and $\dot\eta_r(0)=0$, where $\dot\eta_r$ is the radial derivative (with respect to $\abs{\cdot}$),
%%%\begin{equation}\label{eq:invertmats}
%%%	B(\del_x\varphi_z, \del_\eta\varphi_z)(\hat z,\hat x,\hat \eta)(\dot x(0),\dot\eta(0)) = (\dot x(0),\dot\eta(0))\,.
%%%\end{equation}
%%%Furthermore, 
	the matrix
	\begin{equation}\label{eq:fullyinvmat}
		\begin{pmatrix}
			\del_x\varphi_z & \del_\eta\varphi_z \\
			\del_x\varphi_\eta & \del_\eta\varphi_\eta
		\end{pmatrix}(\hat z,\hat x,\hat\eta) \colon \RR^{N'+n} \to \RR^{N+n}
	\end{equation}
	is injective.
\end{lemma}

\begin{proof}
	%Let $d\pi_L\vert_{(\hat z, \hat \zeta, \hat x, \hat \xi)}$ be injective. 
	Letting $(z(t), x(t),\eta(t))$ be a curve on $\Sigma_\varphi$ with initial condition $(z(0),x(0),\eta(0)) = (\hat z, \hat x, \hat \eta)$, $d\pi_L$ being injective means that
	\begin{align*}
		&{}(\dot z(0), \del_z\varphi_{z}(\hat z, \hat x, \hat \eta)\dot z(0)+ \del_x\varphi_{z}(\hat z, \hat x, \hat \eta)\dot x(0) +  \del_\eta\varphi_{z}(\hat z, \hat x, \hat \eta)\dot\eta(0)) = 0 \\
		&\qquad\implies (\dot x(0), \del_z\varphi_{x}(\hat z, \hat x, \hat \eta)\dot z(0)+ \del_x\varphi_{x}(\hat z, \hat x, \hat \eta)\dot x(0) +  \del_\eta\varphi_{x}(\hat z, \hat x, \hat \eta)\dot\eta(0)) = 0\,,
	\end{align*}
	which is equivalent to 
	\[
(\dot z(0), \del_x\varphi_{z}(\hat z, \hat x, \hat \eta)\dot x(0) +  \del_\eta\varphi_{z}(\hat z, \hat x, \hat \eta)\dot\eta(0)) = 0 \implies (\dot x(0), \del_\eta\varphi_{x}(\hat z, \hat x, \hat \eta)\dot\eta(0)) = 0\,,
	\]
	which concludes the proof of \cref{eq:injmatrix}.

%\begin{equation}\label{eq:injmatrixwith4}
%(\dot z(0), \del_x\varphi_{z}(\hat z, \hat x, \hat \eta)\dot x(0) +  \del_\eta\varphi_{z}(\hat z, \hat x, \hat \eta)\dot\eta(0)) = 0 \implies  = 0\,.
%%\varphi_\eta(\hat z, \hat x, \hat \eta) = 0 \text{ implies that the } (N\times (n+N'))\text{-matrix}\quad \left((\del_\eta\varphi_{z})(\hat z,\hat x,\hat \eta), (\del_x\varphi_{z})(\hat z,\hat x,\hat \eta)\right) \quad\text{is injective}\,.
%\end{equation} 
	We turn to proving the second part. Notice first that by the non-degeneracy of $\varphi$,
	\[
		\begin{pmatrix}\del_\eta\varphi_z \\ \del_\eta\varphi_x \\ \del_\eta\varphi_\eta \end{pmatrix}(\hat z,\hat x,\hat\eta)\cdot \dot \eta(0) = 0 \implies \dot\eta(0)=0\,.
	\]
	Second, since $(z,x,\eta)(t)$ is a curve on $\Sigma_\varphi$, we have
	\[
		0 = \del_t(\varphi_\eta((z,x,\eta)(t))) = \del_z \varphi_\eta(\hat z,\hat x, \hat \eta) \dot z(0) + \del_x\varphi_\eta(\hat z,\hat x, \hat \eta)\dot x(0) + \del_\eta \varphi_\eta(\hat z,\hat x, \hat \eta) \dot \eta(0)\,,
	\]
	so that
	\[
		\dot z(0)=0\,,\dot x(0) = 0 \implies \del_\eta \varphi_\eta(\hat z, \hat x, \hat \eta) \dot \eta(0) = 0\,.
	\]
	Therefore, $\varphi$ being non-degenerate implies 
	\begin{equation}\label{eq:nondegenjustcheckxz}
		\dot z(0)=0\,, \dot x(0)=0\,, \begin{pmatrix}\del_\eta\varphi_z \\ \del_\eta\varphi_x\end{pmatrix}(\hat z,\hat x, \hat\eta)\dot \eta(0) = 0 \implies \dot\eta(0)=0\,.
	\end{equation}

	Now by assumption~\ref{as:n2} we have
	\[
		\dot z(0)=0, \dot x(0)=0,\del_\eta \varphi_x(\hat z, \hat x, \hat \eta) \dot \eta(0) =0 \implies \del_\eta\varphi_z(\hat z,\hat x,\hat \eta)\dot \eta(0) = 0\,,%\dot\eta(0)=0\,.
	\]
	which combined with \cref{eq:nondegenjustcheckxz} gives
	\[
		\dot z(0)=0, \dot x(0)=0,\del_\eta \varphi_x(\hat z, \hat x, \hat \eta) \dot \eta(0) =0 \implies \dot\eta(0)=0\,.
	\]
	Thus, using \cref{eq:injmatrix}, we get
	\[
		(\dot z(0), %\del_z\varphi_{z}(\hat z, \hat x, \hat \eta)\dot z(0)+ 
		\del_x\varphi_{z}(\hat z, \hat x, \hat \eta)\dot x(0) +  \del_\eta\varphi_{z}(\hat z, \hat x, \hat \eta)\dot\eta(0)) = 0 \implies (\dot x(0), \dot\eta(0)) =0\,,
	\]
	so that the matrix $(\del_x\varphi_z, \del_\eta\varphi_z)(\hat z, \hat x, \hat \eta)$ is injective on the vector space
	\[
		W\coloneqq \left\{\begin{pmatrix}\dot x(0)\\ \dot\eta(0)\end{pmatrix} \colon (z,x,\eta)(t) \text{ curve on } \Sigma_\varphi\text{ with } \dot z(0)=0, (z,x,\eta)(0)=(\hat z,\hat x,\hat \eta)\right\}\,.
	\]
	Rewriting $W$, %and using \cite[Prop.~5.38]{zbMATH06034615} 
	we know that
	\[
		W = T_{(\hat x,\hat \eta)}\left(\left\{(x,\eta)\colon (\hat z, x,\eta) \in \Sigma_\varphi\right\}\right) = \mathrm{ker}((\del_x\varphi_\eta,\del_\eta\varphi_\eta)(\hat z,\hat x,\hat\eta))\,. 
	\]
	Locally, of course, $\RR^{N'+n} = W \oplus W^\perp$, and we note that if $w_1,w_2\in W^\perp$ with $(\del_x\varphi_\eta,\del_\eta\varphi_\eta)(\hat z, \hat x,\hat \eta) w_1 = (\del_x\varphi_\eta,\del_\eta\varphi_\eta)(\hat z, \hat x,\hat \eta) w_2$, we have $w_1-w_2 \in W^\perp \cap W$ so that $w_1=w_2$. This is to say that the matrix
	\begin{equation}\label{eq:injmatonperp}
		(\del_x\varphi_\eta,\del_\eta\varphi_\eta)(\hat z, \hat x,\hat \eta)\colon W^\perp \to \RR^n
	\end{equation}
	is injective. 

	Now write any element in $\RR^{N'+n}$ as $w+ w^\perp$ with $w\in W,w^\perp\in W^\perp$. Then
	\[
		\begin{pmatrix}
			\del_x\varphi_z & \del_\eta\varphi_z \\ 
			\del_x\varphi_\eta & \del_\eta\varphi_\eta
			\end{pmatrix}(\hat z,\hat x,\hat\eta) (w+w^\perp) = \begin{pmatrix}(\del_x\varphi_z,\del_\eta\varphi_z)(\hat z,\hat x,\hat\eta)(w+w^\perp)\\ (\del_x\varphi_\eta,\del_\eta\varphi_\eta)(\hat z,\hat x,\hat\eta)w^\perp\end{pmatrix}
	\]
	where the RHS can only be zero if $w^\perp = 0$ by \cref{eq:injmatonperp}, and then $w=0$ because $(\del_x\varphi_z, \del_\eta\varphi_z)(\hat z, \hat x, \hat \eta)$ is injective on $W$. This concludes the proof of \cref{eq:fullyinvmat}.
\end{proof}

\end{document}
